\section{Scope, conventions, and main results}

The source manuscript~\cite{ProveIt2026} assembles a substantial collection of exact
identities, experimental reductions, and open questions about
polylogarithmic constants. Its Gaussian depth chapter contains several
precise conjectural rows for which a proof can be obtained without
assuming a period conjecture or promoting numerical evidence to equality.
We concentrate on those rows, their natural generalizations, and the
reliability of the claimed remaining directions. The resulting article
is a focused continuation of Chapters~4 and~6, with implications for the
computational methodology of the wider manuscript. It is not an assertion
that every claim in all chapters has been exhaustively checked.

\subsection{Results and their logical status}

\begin{center}
\small
\begin{tabular}{@{}P{0.25\textwidth}P{0.34\textwidth}P{0.33\textwidth}@{}}
\toprule
Source or family & Result established here & Scope of the conclusion\\
\midrule
Five Gaussian doubles, weight six & All five displayed rows proved;
two exact derivations & Coefficients unchanged; no independence claim\\[4pt]
Gaussian doubles, even weight & Finite rational coefficient formula;
seven explicit weight-eight rows & Complete reduction of the indicated
imaginary components\\[4pt]
Three Gaussian triples, weight four & All three displayed rows proved
from one uniform identity & Coefficients unchanged\\[4pt]
One $2$ among any number of $1$'s & Classical formula at every position
and every weight & At most one imaginary direction modulo the specified
lower-weight products\\[4pt]
Primitive sixth-root endpoints & Same-argument closure and an explicit
rational parity component & Every weight and every position of the $2$\\[4pt]
Four mixed-root candidates & Explicit Gaussian and Eisenstein reductions
& These candidates do not add directions to the stated single-value basket\\[4pt]
Numerical evaluation & Exact rational H\"older certificates and an
independent replay & Certified enclosures; equality follows from the
analytic proofs\\[4pt]
Local editorial assertions & Corrected color exchange, diagonal value,
and shifted-polygamma parity & Explicit replacement text and source patch\\
\bottomrule
\end{tabular}
\end{center}

There are several different meanings of ``reduction'' in this subject.
An identity may express a constant through a chosen collection of other
constants. A relation matrix may have a specified rank over $\mathbb Q$.
Neither statement proves that the corresponding numerical constants are
independent. An integer-relation search that fails to find a relation
does not prove nonexistence. We use ``basis'' only when discussing a
chosen coordinate basket or a proved vector-space basis, and explicitly
identify which is meant.

\subsection{Series convention and iterated-integral dictionary}

For positive integers $s_1,\ldots,s_d$, our colored convention is
\begin{equation}\label{eq:mpl-definition}
 \Li_{s_1,\ldots,s_d}(z_1,\ldots,z_d)
 =\sum_{n_1>\cdots>n_d\ge1}
 \frac{z_1^{n_1}\cdots z_d^{n_d}}
 {n_1^{s_1}\cdots n_d^{s_d}}.
\end{equation}
The weight is $s_1+\cdots+s_d$ and the series depth is $d$.
The abbreviation $\Li_{s_1,\ldots,s_d}(z)$ means
$(z_1,z_2,\ldots,z_d)=(z,1,\ldots,1)$. In particular,
$\Li_{a,b}(i)$ and $\Li_{a,b}(i,i)$ are different objects.
Three families will be distinguished throughout:
\[
 D_{a,b}(z)=\Li_{a,b}(z,1),\qquad
 U_{a,b}(z)=\Li_{\{1\}^a,2,\{1\}^b}(z),\qquad
 M_{a,b}(x)=\Li_{a,b}(x,x^{-1}).
\]

With the outermost-letter-first integral convention
$G(a_1,\ldots,a_w;z)=\int_0^zG(a_2,\ldots,a_w;t)\,dt/(t-a_1)$,
the exact dictionary is
\begin{equation}\label{eq:mpl-gpl}
 \Li_{s_1,\ldots,s_d}(z_1,\ldots,z_d)
 =(-1)^dG\left(0^{s_1-1},z_1^{-1},
 0^{s_2-1},(z_1z_2)^{-1},\ldots,
 0^{s_d-1},(z_1\cdots z_d)^{-1};1\right).
\end{equation}
For the single-endpoint family, an equivalent positive-form word uses
$dt/t$ and $dt/(1-t)$ with endpoint $z$. The factor $(-1)^d$ in
\eqref{eq:mpl-gpl} accounts for the convention difference. This
dictionary is also the input specification of the rational evaluator.
Section~\ref{sec:certificates} proves that the convergent boundary series
used here agree with the ordinary iterated integrals; no divergent
regularization is hidden in the notation.

We write $\ell=\log2$, $G=\beta(2)$, and
\[
 \beta(s)=\sum_{n\ge0}\frac{(-1)^n}{(2n+1)^s},\qquad
 h=\frac{1+i}{2},\qquad
 \mu_s=\Re\Li_s(h),\qquad\lambda_s=\Im\Li_s(h).
\]
Locally in Section~\ref{one2:sec:general}, $w$ denotes the transformed
argument $(1-z)^{-1}$; elsewhere a weight is explicitly identified as
such. All logarithms and classical polylogarithms have their principal
branches unless a continuation path is specified.

\subsection{Relation to established results}

Panzer's parity theorem already gives functional reductions of multiple
polylogarithms and explicit formulas in depths two and three
\cite{Panzer2017}. The present double identities specialize its depth-two
formula, with an explicit singular-limit calculation and an independent
differential verification. Umezawa subsequently gave a general explicit
parity formula \cite{Umezawa2025}; accordingly, this article does not claim
the first general or all-depth parity theorem.

The classical reducibility of the height-one Nielsen family is also
established. Borwein and Straub's complementary-argument formula
\cite[Eq.~(33)]{BorweinStraub2015} is closely related to the logarithmic
moment used below. Our useful form has argument $(1-z)^{-1}$, explicit
coefficients for every position of the $2$, carefully fixed branches and
endpoint constants, and a direct application to the three manuscript
rows. Similarly, H\"older convolution and its use in multiple
polylogarithm evaluation are classical \cite{BBBL2001,VollingaWeinzierl2005}.
We supply an explicit rational certificate, a uniform precision bound,
and a replayable implementation for the relevant Gaussian letters.

These boundaries matter to the interpretation of the results. The
contribution consists of proof upgrades, explicit extensions of the
pinned tables, independently checkable formulas, and corrections to
specific claims. Global bibliographic novelty of each isolated numerical
identity is not asserted. The four mixed-root reductions are also
presented as corroborated corrections; related companion work in this
research program has already reported these reductions.
